% Segment Tree with Lazy Propagation

Range update and range query in $O(\log N)$.
\begin{lstlisting}
struct SegTreeLazy {
    int n; vector<int> tree, lazy;
    SegTreeLazy(int n) : n(n), tree(4 * n, 0), lazy(4 * n, 0) {}
    void push(int node, int l, int r) {
        if (!lazy[node]) return;
        int mid = (l + r) >> 1;
        tree[node << 1] += (mid - l + 1) * lazy[node];
        lazy[node << 1] += lazy[node];
        tree[node << 1 | 1] += (r - mid) * lazy[node];
        lazy[node << 1 | 1] += lazy[node];
        lazy[node] = 0;
    }
    void update(int node, int l, int r, int ql, int qr, int val) {
        if (ql <= l && r <= qr) {
            tree[node] += (r - l + 1) * val;
            lazy[node] += val;
            return;
        }
        push(node, l, r);
        int mid = (l + r) >> 1;
        if (ql <= mid) update(node << 1, l, mid, ql, qr, val);
        if (qr > mid) update(node << 1 | 1, mid + 1, r, ql, qr, val);
        tree[node] = tree[node << 1] + tree[node << 1 | 1];
    }
    int query(int node, int l, int r, int ql, int qr) {
        if (ql <= l && r <= qr) return tree[node];
        push(node, l, r);
        int mid = (l + r) >> 1, res = 0;
        if (ql <= mid) res += query(node << 1, l, mid, ql, qr);
        if (qr > mid) res += query(node << 1 | 1, mid + 1, r, ql, qr);
        return res;
    }
};
\end{lstlisting}
